\section{Vektorrum}
Alle vektorrum skal opfylde aksiomerne:
\begin{itemize}
\itemsep0em 
    \item \( \mathbf{u} + (\mathbf{v} + \mathbf{w}) = (\mathbf{u} + \mathbf{v}) + \mathbf{w} \)
    \item \( \mathbf{u} + \mathbf{v} = \mathbf{v} + \mathbf{u} \)
    \item Der findes en nulvektor \( \mathbf{0} \) sådan at \( \mathbf{u} + \mathbf{0} = \mathbf{u} \)
    \item For hver vektor \( \mathbf{u} \) findes en modsat vektor \( -\mathbf{u} \) sådan at \( \mathbf{u} + (-\mathbf{u}) = \mathbf{0} \)
    \item \( c \cdot (d \cdot \mathbf{u}) = (c \cdot d) \cdot \mathbf{u} \)
    \item \( 1 \cdot \mathbf{u} = \mathbf{u} \)
    \item \( c \cdot (\mathbf{u} + \mathbf{v}) = c \cdot \mathbf{u} + c \cdot \mathbf{v} \)
    \item \( (c + d) \cdot \mathbf{u} = c \cdot \mathbf{u} + d \cdot \mathbf{u} \)
\end{itemize}
\kilde{Definition 10.1.1}

\subsection{Standardbasisvektoren}
Standardbasisvektorerne er de lineært uafhængige vektorer i et vektorrum, som har et 1 i en position og 0 i alle andre positioner.
I \(\mathbb{R}^3\) er standardbasisvektorerne:
\[\mathbf{e}_1 = \begin{bmatrix} 1 \\ 0 \\ 0 \end{bmatrix}, \quad \mathbf{e}_2 = \begin{bmatrix} 0 \\ 1 \\ 0 \end{bmatrix}, \quad \mathbf{e}_3 = \begin{bmatrix} 0 \\ 0 \\ 1 \end{bmatrix}\]


\subsection{Find Beta-koordianter for en vektor ved række-reduktion}
\label{sec:beta-koordinater}
Givet vektor $\mathbf{v}$ og basis $\beta = \{\mathbf{b}_1, \mathbf{b}_2, \ldots, \mathbf{b}_n\}$ for vektorrum $V$, kan vi finde $[\mathbf{v}]_\beta$ ved at række-reducere \([\mathbf{b}|\mathbf{v}]\) matricen til reduceret trappeform.
\[
\begin{bmatrix}
    \vert & & \vert & \vert \\
    \mathbf{b}_1 & \cdots & \mathbf{b}_n & \mathbf{v} \\
    \vert & & \vert & \vert
\end{bmatrix}
\xrightarrow{\text{Række-operationer}}
\begin{bmatrix}
    1 & \cdots & 0 & \vert \\
    \vdots & 1 & \vdots & \mathbf{v}_\beta \\
    0 & \cdots & 1 & \vert
\end{bmatrix}
\]

\subsection{Om beta-koordinater}
\[ [\mathbf{v} + \mathbf{u}]_\beta = [\mathbf{v}]_\beta + [\mathbf{u}]_\beta \]
\[ [c \cdot \mathbf{v}]_\beta = c \cdot [\mathbf{v}]_\beta \]
\kilde{Lemma 10.2.3}

\[ c_1 \cdot \mathbf{v}_1 + c_2 \cdot \mathbf{v}_2 + \cdots + c_n \cdot \mathbf{v}_n = \mathbf{0} \iff c_1 \cdot [\mathbf{v}_1]_\beta + c_2 \cdot [\mathbf{v}_2]_\beta + \cdots + c_n \cdot [\mathbf{v}_n]_\beta = \mathbf{0} \]
\kilde{Sætning 10.2.4}

\subsection{Ordnet basis af vektorrum}
For vektorrum \(V=\spanop(v_1,v_2,\cdots,v_n)\) kan den ordnede basis findes som beskrevet i sektion \ref{sec:ordnet-basis-udspaending}

\subsubsection{Er givent sæt vektorer en ordnet basis?}
\begin{enumerate}
    \item Der skal være lige så mange vektorer som dimensionen af vektorrummet.
    \item Vektorerne skal være lineært uafhængige.
\end{enumerate}
\kilde{Sætning 10.2.7}

De er lineært uafhængige hvis matricen dannet af vektorerne har pivot i alle søjler efter række-reduktion til trappeform.

\subsection{Bestem basisskiftematrix}
\subsubsection{Ved individuelle søjler}
Givet to basiser
$\gamma = \{\mathbf{w}_1, \mathbf{w}_2, \cdots, \mathbf{w}_n\}$ og $\beta = \{\mathbf{v}_1, \mathbf{v}_2, \cdots, \mathbf{v}_n\}$
for vektorrum $V$, kan vi bestemme basisskiftematricen $_\beta M _\gamma$ (Fra $\gamma$ til $\beta$) 
ved at finde koordinaterne for hver eneste basisvektor i $\gamma$ i forhold til basisen $\beta$ som i sektion \ref{sec:beta-koordinater}

De resulterende koordinatvektorer udgør søjlerne i basisskiftematricen:
\[_\beta M_\gamma = \begin{bmatrix} [\mathbf{w}_1]_\beta & [\mathbf{w}_2]_\beta & \cdots & [\mathbf{w}_n]_\beta \end{bmatrix} \kildetag{Sætning 10.3.1} \]

\subsubsection{Basisskiftematricer skifter basis}
\[_\gamma M_\beta \cdot [\mathbf{v}]_\beta = [\mathbf{v}]_\gamma \kildetag{Sætning 10.3.1}\]

\subsubsection{Sætilfælde: Fra den ordnede standarbasis}
Enhver basisskiftematix $_\beta M _\epsilon$ (Fra standardbasis til $\beta$) kan findes ved at danne en matrix med basisvektorerne i $\beta$ som søjler:
\[_\beta M_\epsilon = \begin{bmatrix} \mathbf{v}_1 & \mathbf{v}_2 & \cdots & \mathbf{v}_n \end{bmatrix}\]
Det er altså bare at skrive basisvektorerne ved siden af hinanden som søjler i en matrix.

\subsubsection{Hvis du kender den modsattet basisskiftematix}
Hvis du kender basisskiftematricen $_\beta M _\gamma$ (Fra $\gamma$ til $\beta$), kan du finde den modsatte basisskiftematix $_\gamma M _\beta$ (Fra $\beta$ til $\gamma$) ved at tage den inverse af $_\beta M _\gamma$:
\[_\gamma M _\beta = \left(_\beta M _\gamma\right)^{-1} \kildetag{Lemma 10.3.2}\]

\subsubsection{Kombination af basisikftematricer}
Hvis du kender basisskiftematricerne $_\beta M _\gamma$ (Fra $\gamma$ til $\beta$) og $_\gamma M _\delta$ (Fra $\delta$ til $\gamma$), kan du finde basisskiftematricen $_\beta M _\delta$ (Fra $\delta$ til $\beta$) ved at multiplicere de to matricer:
\[_\beta M _\delta = _\beta M _\gamma \cdot _\gamma M _\delta \kildetag{Lemma 10.3.2}\]
\bemærk{Rækkefølgen af multiplikationen er vigtig}

\subsection{Underrum}
For underrum $W$ af vektorrum $V$ gælder $\dim(W) \leq \dim(V)$

\kilde{Lemma 10.4.1}

\

For 2 vektorer $\mathbf{u}, \mathbf{v} \in W$ gælder $\mathbf{u} + c \cdot \mathbf{v} \in W$ for alle skalarer $c \in \mathbb{F}$.

\kilde{Lemma 10.4.2}

\eksempel{
    er mængden $\{ a + bz \mid a, b \in \mathbb{C}\}$ et underrum af vektorrummet af alle komplekse polynomier $\mathbb{C}[z]$?
    \[
    (a+bz) + c \cdot (k+jz) = a + ck + bz + cjz = (a + ck) + (b + cj)z
    \]
    Det er sandt, da det kunne omskrives til at være et polynomium af samme form.
}

\subsubsection{Find dimensionen}
Dimensionen af et underrum er antallet af variabler der kan vælges frit.

\eksempel{
    \[
    \begin{bmatrix}
        a_{11} & a_{12} & a_{13} \\
        0 & a_{22} & a_{23} \\
        0 & 0 & a_{33} \\
    \end{bmatrix}
    \]

    Dimensionen af alle \(\mathbb{R}^{3 \times 3}\) øvre-trekantsmatricer er $3 + 2 + 1 = \frac{3 \cdot (3+1)}{2} = 6$
}

\subsection{Ordnet basis for underrum}
Vektorer er en ordnet basis for et underrum hvis de er:

Lineært uafhængige, i underrummet, og der er samme antal vektorer som dimensionen af underrummet.

Se sektion \ref{sec:ordnet-basis-udspaending} for ordnede baser.

%Givet en ordnet basis $\beta$ for $V$, og vektorer der udgør et undderum af $V$ $v_1, v_2, \cdots, v_n$, kan du finde en ordnet basis for $\spanop(v_1, v_2, \cdots, v_m)$ ved finde pivotsøjlerne af $[v_1]_\beta, [v_2]_\beta, \cdots, [v_m]_\beta$. Da er de originale den ordnede basis.
%\kilde{Sætning 10.4.4}